\begin{pastexam}[label={px:nchu-115-is-20}]{中興資工所乙組 115 學年度 第 20 題}
（單選；每題 2 分，答錯不倒扣）

You want to determine whether two large files (gigabytes in size) are identical without reading
the entire files unless absolutely necessary. What approach is best?
\begin{enumerate}[label=\Alph*.,nosep]
\item Compute and compare cryptographic hash values (e.g., SHA-256) of both files.
\item Load both files fully into memory and compare them byte by byte.
\item Randomly guess whether the files are identical and confirm only if the guess seems wrong.
\item Run insertion sort on the bytes of each file and compare the sorted results.
\end{enumerate}
\end{pastexam}

\begin{pastsoln}
答案是 A（出題者的本意）。

先說題目的問題：算 SHA-256 一樣要把兩個檔案從頭讀到尾，做不到題目說的「除非必要，不讀完
整個檔案」。 真正少讀的做法是先比檔案大小（不同就一定不同），再一段一段讀進來比，遇到第一個
不同的位元組就停；這不在選項裡。 雜湊的好處在別的情況：檔案在兩台機器上時只要傳 $32$ 位元組
的雜湊值，或同一個檔案要和很多檔案比時，雜湊值算一次存起來就能重複用。 四個選項裡只有 A
合理，照出題者的意思選 A。
\begin{itemize}
\item A：兩個檔案相同，雜湊值一定相同；雜湊值不同，檔案一定不同。 雜湊值相同而檔案不同是碰撞，
雜湊值有 $2^{256}$ 種，密碼學雜湊函數又設計成很難找到碰撞，實務上可以當成相同。 和資結書
第~12.1~節的雜湊函數相比，這裡的雜湊值不拿來當陣列索引，而是當成檔案的「指紋」。
\item B 錯：GB 等級的檔案整個讀進記憶體可能放不下；而且逐位元組比也不需要全部載入，邊讀邊比
就好。 這個選項做的事比必要的還多。
\item C 錯：用猜的沒有任何保證。
\item D 錯：排序後相同只表示兩個檔案的位元組「組成」相同，順序不同的兩個檔案也會一樣，判斷是
錯的；而且對幾十億個位元組做插入排序是 $\Theta(n^2)$。
\end{itemize}
\end{pastsoln}
